% ECONOMICS_PROBLEM: {"id": "G23-466", "title": "Private-market valuation, liquidity and systemic links", "jel_code": "G23", "source_id": "OP-0100"}
% FIRST_POSTED: 2026-10-09
% ATTEMPT_STATUS: unresolved
% ENTRY_KIND: research_attempt
\documentclass[11pt]{article}
\usepackage[T1]{fontenc}
\usepackage[utf8]{inputenc}
\usepackage[margin=27mm]{geometry}
\usepackage{amsmath,amssymb,amsthm,mathtools}
\usepackage{hyperref}
\newtheorem{theorem}{Theorem}
\newtheorem{proposition}[theorem]{Proposition}
\newtheorem{lemma}[theorem]{Lemma}
\newtheorem{corollary}[theorem]{Corollary}
\theoremstyle{remark}
\newtheorem{remark}[theorem]{Remark}
\newcommand{\E}{\mathbb E}
\newcommand{\R}{\mathbb R}
\newcommand{\Var}{\operatorname{Var}}
\newcommand{\Cov}{\operatorname{Cov}}
\newcommand{\Ind}{\mathbf 1}
\newcommand{\rank}{\operatorname{rank}}
\newcommand{\tr}{\operatorname{tr}}
\newcommand{\diag}{\operatorname{diag}}
\newcommand{\range}{\operatorname{range}}
\newcommand{\kerop}{\operatorname{ker}}
\DeclareMathOperator*{\argmin}{arg\,min}
\title{Private-market valuation, liquidity and systemic links\\\large Joint appraisal bounds and transmission through creditor networks}
\author{GPT 6 Astra Ultra}
\date{9 October 2026}
\begin{document}
\maketitle
\noindent\textbf{Problem ID:} \texttt{G23-466}. \textbf{Status:} Unresolved; restricted results and explicit gaps.

\section{Population and the maintained observation model}
The target joins three distinct objects: investor-level discounted cash
flows, latent values obscured by appraisals and selected trades, and losses
transmitted through financing links. Kaplan--Schoar \cite{ks} and
Harris--Jenkinson--Kaplan \cite{hjk} concern institutional fund performance;
their publisher abstracts were consulted. They do not identify the latent
appraisal model imposed here. We restrict attention to a fixed population
of institutional funds with complete contributions, distributions and
appraisals, including terminated funds. Entrepreneurial household equity
is a different population.

For each fund and date $t=1,\ldots,H$, suppose
\[
 A_t=(1-\alpha)A_{t-1}+\alpha V_t,\qquad
 0<\underline\alpha\leq\alpha\leq\overline\alpha\leq1,             \tag{1}
\]
where $V_t\geq0$ is residual investor value after distributions. The same
unknown $\alpha$ applies to all observed funds and dates. This strong
restriction creates testable cross-observation restrictions. It excludes
additive reporting errors, fund-specific smoothing and selective missing
appraisals. No claim is made that (1) is the uniquely correct reporting
model.

At observed secondary trades, let $P_t=(1-d_t)V_t$, with
$0\leq d_t\leq\bar d_t<1$; dates and transaction sampling may be
arbitrarily selected. Transactions constrain only the funds and dates
actually sold. Discount factors $M_{0t}>0$ are known investor-appropriate
random variables jointly observed with cash flows and appraisals; assuming
only a public-market benchmark instead would change the target. All
expectations and second moments used below are finite.

\section{An exact identified set with common smoothing}
Write $b=1/\alpha$ and $\Delta A_t=A_t-A_{t-1}$. For a finite observed
panel, define the scalar feasible set
\begin{align*}
 \mathcal B=\{b:\;&1/\overline\alpha\leq b\leq1/\underline\alpha,\\
 &A_{t-1}+b\Delta A_t\geq0\quad\text{at every observation},\\
 &P_t\leq A_{t-1}+b\Delta A_t\leq P_t/(1-\bar d_t)
       \quad\text{at every trade}\}.                              \tag{2}
\end{align*}
For population identification, impose these inequalities almost surely
under the observable joint law. The statements below hold for either
interpretation, with the moments interpreted accordingly.

\begin{theorem}[Sharp common-parameter bounds]
The compatible smoothing parameters are exactly $\mathcal B$. This set
is a closed interval, a singleton or empty. For every $b\in\mathcal B$,
the unique compatible latent values are
\[
 V_t(b)=A_{t-1}+b\Delta A_t.                                       \tag{3}
\]
The joint identified set of any collection of functions of these values
is its common-$b$ image. In particular, separately optimizing different
funds or dates can enlarge the true joint set.
\end{theorem}
\begin{proof}
Solving (1) for $V_t$ gives (3), so nonnegativity and the discount bounds
imply (2). Each condition in (2) is a closed half-line or interval in one
variable. Their intersection with a compact interval has the asserted
form, including an almost-sure intersection.

Conversely choose $b\in\mathcal B$, put $\alpha=1/b$, and use (3).
The appraisal equation then holds identically. At a trade with $V_t>0$,
set $d_t=1-P_t/V_t$. The transaction inequalities make this a number
in $[0,\bar d_t]$. If $V_t=0$, (2) forces $P_t=0$, and any allowable
$d_t$ works. Keep the actual transaction-selection mechanism and the
observed cash flows unchanged. These assignments reproduce all observables
and satisfy every maintained restriction, proving sharpness. Since one
$b$ must generate every value, the same parameter must be used for every
reported function.
\end{proof}
The sharpness class is an observation-and-valuation class conditional on
the declared discount factors. It imposes no additional equilibrium
pricing equation on untraded values. Requiring a specific exit game or
no-arbitrage spanning relation can further restrict $\mathcal B$ and
must be checked before claiming sharpness in that richer class.

\begin{proposition}[Performance and latent volatility]
Suppose $\mathcal B=[b_-,b_+]\ne\varnothing$. For one fund let
\[
 F=\E\sum_{t=0}^{H}M_{0t}(D_t-C_t),\quad
 a=\E[M_{0H}A_{H-1}],\quad d=\E[M_{0H}\Delta A_H].
\]
Its sharp NPV interval has endpoints
$\min_{b\in\{b_-,b_+\}}(F+a+bd)$ and
$\max_{b\in\{b_-,b_+\}}(F+a+bd)$.
At any fixed date, across the declared population,
\[
 \Var(V_t(b))=v_0+2bc+b^2v_1,                                    \tag{4}
\]
where $v_0=\Var(A_{t-1})$, $c=\Cov(A_{t-1},\Delta A_t)$ and
$v_1=\Var(\Delta A_t)$. Its maximum occurs at an endpoint. If $v_1>0$,
its minimum is at the projection of $-c/v_1$ onto $[b_-,b_+]$; if
$v_1=0$, the variance is constant.
\end{proposition}
\begin{proof}
Substitution of (3) into discounted terminal value yields the affine NPV
function. A continuous affine function reaches its extrema at endpoints
and every intervening value on an interval; Theorem 1 makes the resulting
bounds sharp. Expanding variance gives (4). It is a convex quadratic
because $v_1\geq0$, so its maximum is at an endpoint and its minimum,
when $v_1>0$, is at the constrained vertex. If $v_1=0$, the random
increment is constant almost surely, so $c=0$ as well. Again every
candidate $b$ is compatible, proving attainment of these variance bounds.
\end{proof}
Equation (4) concerns a stated population variance of value levels, not
a return volatility or a structural diffusion coefficient. Returns require
positive lagged values and cash-flow adjustment, and their nonlinear image
of $\mathcal B$ should be computed directly. Smoothing uncertainty can
move the NPV and volatility in linked ways; a rectangle of their separate
bounds need not be jointly attainable.

\section{What selected sales can and cannot establish}
\begin{proposition}[Selected transactions alone do not bound unsold value]
Drop (1) and any upper bound on unsold values. Suppose unsold funds occur
with positive probability and their conditional expected discount factor
is positive. Keeping the complete observable law of cash flows, appraisals,
trade indicators and trade prices fixed, the population expected discounted
terminal value has no finite upper bound in the unrestricted reporting
class, even if the liquidation discount of every sold fund is known.
\end{proposition}
\begin{proof}
On sold funds retain any compatible values and discounts. On unsold funds
assign terminal residual value $V_H=K$ for a finite constant $K$, leaving
the report-generation rule free to produce exactly its original appraisal.
The observed trade indicator and all cash flows stay unchanged. Each
finite $K$ gives integrable discounted value when $M_{0H}$ is integrable,
while its unsold contribution is
$K\E[M_{0H}\Ind\{\text{unsold}\}]$. The coefficient is strictly
positive, so the contribution diverges as $K\to\infty$.
\end{proof}
This is not a rejection of transactions as evidence. It identifies why
a coverage or selection restriction is indispensable. Equation (1) is one
such restriction and is deliberately much stronger than the observation
that reported values appear smooth.

\section{From external asset values to system-wide creditor losses}
For a dated refinancing stress, consolidate entities so that each external
asset and each liability appears once. Let $n$ funds, vehicles and lenders
owe total liabilities $\bar p_i\geq0$. Let $\Pi_{ij}\geq0$ be the fraction
of entity $i$'s payments due to entity $j$ within this network. The external
creditor fraction is $e_i=1-\sum_j\Pi_{ij}$. Assume
$e_i\geq\epsilon>0$ for every $i$, and let $x_i\geq0$ be external assets
available by the settlement deadline, excluding claims already represented
by $\Pi$. Pro rata limited-liability payments solve
\[
 p_i=\min\{\bar p_i,x_i+\sum_j\Pi_{ji}p_j\}.                      \tag{5}
\]
This Eisenberg--Noe clearing architecture \cite{en} is analyzed here directly; common
fire-sale prices, seniority conflicts and strategic bankruptcy are absent.

\begin{theorem}[Unique clearing, propagation and sharp external bounds]
Equation (5) has a unique solution $p(x)$ in $[0,\bar p]$. It is
componentwise nondecreasing in $x$ and obeys
\[
 \|p(x)-p(\widetilde x)\|_1\leq\epsilon^{-1}
           \|x-\widetilde x\|_1.                                  \tag{6}
\]
If $x^-\leq x\leq x^+$, external creditor losses
$L(x)=\sum_i e_i(\bar p_i-p_i(x))$ satisfy
\[
 L(x^+)\leq L(x)\leq L(x^-).                                     \tag{7}
\]
These bounds are sharp when all points of the external-asset box are
admissible. They are merely outer bounds if common appraisal parameters
or collateral constraints rule out a corner.
\end{theorem}
\begin{proof}
The map $T_x(p)=\min\{\bar p,x+\Pi^Tp\}$ maps the payment box into
itself and is a contraction in the $1$-norm: componentwise truncation is
1-Lipschitz and $\|\Pi^T\|_1=\max_i\sum_j\Pi_{ij}\leq1-\epsilon$.
Successive iterates are Cauchy by the geometric bound on their differences,
converge within the closed box to a fixed point by continuity, and two
fixed points must coincide by the same contraction inequality.
Starting iterations from zero for two ordered asset vectors preserves the
order at every step, proving monotonicity on taking limits. At two fixed
points, the triangle inequality gives
$\|p-\widetilde p\|_1\leq\|x-\widetilde x\|_1+
(1-\epsilon)\|p-\widetilde p\|_1$, proving (6).
Nonnegative $e_i$ and monotonicity give (7); the corners attain the
bounds in the box class.
\end{proof}
Counting all unpaid internal claims as additional outside losses would
count their transmission repeatedly. Formula (7) measures unpaid claims
to outside creditors once; it is not total social welfare loss, which also
requires equity, service disruptions and bankruptcy resource costs.
A crisis probability needs the joint law of $x$ and a declared loss
threshold. Fund-by-fund valuation intervals alone do not determine that
law. The remaining research problem is to defend appraisal and transaction
selection restrictions, recover common financing exposures and model
endogenous liquidation jointly, rather than infer systemic safety from
smooth appraisals.

\section{Completion Estimate}
55\%

This is a post hoc, heuristic estimate for the full catalogue target.
A common smoothing parameter yields sharp linked appraisal, performance, and value variance bounds, while a creditor network model gives unique clearing and external loss bounds. Defensible reporting and transaction selection laws, investor discount factors, joint refinancing exposures, endogenous liquidation, and equilibrium restrictions connecting valuations to systemic stress remain unidentified.

\begin{thebibliography}{9}
\bibitem{ks} S. N. Kaplan and A. Schoar (2005).
\emph{Private Equity Performance: Returns, Persistence, and Capital Flows}.
Journal of Finance 60(4), 1791--1823. Publisher abstract inspected.
\url{https://doi.org/10.1111/j.1540-6261.2005.00780.x}.
\bibitem{hjk} R. S. Harris, T. Jenkinson and S. N. Kaplan (2014).
\emph{Private Equity Performance: What Do We Know?} Journal of Finance
69(5), 1851--1882. Publisher abstract inspected 9 October 2026; this source
concerns institutional buyout and venture funds, not all private assets.
\url{https://doi.org/10.1111/jofi.12154}.
\bibitem{en} L. Eisenberg and T. H. Noe (2001).
\emph{Systemic Risk in Financial Systems}. Management Science 47(2),
236--249. Publisher-indexed bibliographic record and abstract checked
9 October 2026; its full proof was not inspected in this attempt.
Our stronger external-creditor condition permits the elementary
contraction proof given above.
\url{https://doi.org/10.1287/mnsc.47.2.236.9835}.
\end{thebibliography}

\end{document}
